\documentclass[tikz,border=3pt]{standalone}
\input{tikz-tcee}
\begin{document}
\hyphenpenalty=10000 \exhyphenpenalty=10000
% Efecto bola de nieve: tipo implícito frente al crecimiento nominal
\begin{tikzpicture}[tcee, font=\footnotesize,
  k/.style={draw=#1, fill=#1!12, rounded corners=2pt, line width=0.7pt, text width=3.75cm, align=center, minimum height=1.25cm}]
  \node[k=dicfbcf] at (0,0) {\textbf{$i<\gamma^N$}\\ el crecimiento reduce la ratio por sí solo: basta un \textbf{déficit primario} pequeño};
  \node[k=dicparo] at (4.1,0) {\textbf{$i>\gamma^N$}\\ hacen falta \textbf{superávits primarios}, mayores cuanta más deuda};
\end{tikzpicture}
\end{document}
